\section{TODD：早期棄却法の改良}

以前から提案していたものの，実装上の問題により実験を一時停止していた
早期棄却法について，実装を修正し，再評価を行った．
本実装では，従来の基底再利用法を土台とし，
完全な $\chi$ 行列に対する判定処理を行う前に，
小規模な部分行列を用いた低コストな判定を挿入する．
これにより，削減不可能であると安全に判断できる候補を早い段階で棄却し，
後続の高コストな基底更新や零空間計算を省略することを目的とする．

前回から，提案手法の各機能
（$\chi$ 行列の圧縮，基底再利用，$\chi$ 行列の差分更新，早期棄却）
が実行時間に与える影響を評価するため，
処理の統計情報を記録する機能を実装している．
具体的には，候補数，早期棄却数，基底の再構築回数，
局所更新回数，零空間計算回数，および各処理時間を記録できるようにした．
ただし，本報告ではまず早期棄却法の有効性を確認することを主目的とし，
各統計量に基づく詳細な要因分析は今後の課題とする．

\subsection{手法の説明}

TODD における各候補ペア $(c_1,c_2)$ に対して，
通常は候補に対応する $\chi$ 行列を用いて削減可能性を判定する．
この判定では，零空間ベクトル $y$ が存在し，
\[
    y[c_1] \oplus y[c_2] = 1
\]
を満たすかを調べる必要がある．
しかし，完全な $\chi$ 行列に対する基底更新や零空間計算は高コストであり，
多数の候補ペアに対してこの処理を行うことが実行時間の主要因となる．

本手法では，完全な $\chi$ 行列に対する処理を行う前に，
その一部の行をランダムに取り出した小型 $\chi$ 行列を構成し，
この小型行列に対して行空間への所属判定を行う．
完全な $\chi$ 行列を $M$，その行部分集合からなる小型 $\chi$ 行列を $M'$ とする．
$M'$ は $M$ の行の一部のみからなるため，
\[
    \operatorname{Row}(M') \subseteq \operatorname{Row}(M)
\]
が成り立つ．

候補ペア $(c_1,c_2)$ に対して，
$e_{c_1}+e_{c_2}$ を，$c_1$ 番目と $c_2$ 番目のみが $1$ であるベクトルとする．
このとき，
\[
    e_{c_1}+e_{c_2} \in \operatorname{Row}(M')
\]
が成り立つならば，
\[
    e_{c_1}+e_{c_2} \in \operatorname{Row}(M)
\]
も成り立つ．
線形代数の基本性質より，任意の
$y \in \operatorname{Null}(M)$ に対して，
$\operatorname{Row}(M)$ に属するベクトルとの内積は $0$ となる．
したがって，
\[
    (e_{c_1}+e_{c_2}) \cdot y = 0
\]
であり，
\[
    y[c_1] \oplus y[c_2] = 1
\]
を満たす零空間ベクトル $y$ は存在しない．
よって，この候補ペアは削減不可能であると判断できる．

以上より，小型 $\chi$ 行列 $M'$ において
$e_{c_1}+e_{c_2}$ が行空間に含まれる場合，
その候補は安全に棄却できる．
この場合，完全な $\chi$ 行列に対する状態更新，行空間基底の更新，
および零空間計算を行わずに次の候補へ進む．
一方，小型行列で棄却できなかった候補については，
従来の基底再利用法と同様に，完全な $\chi$ 行列の状態および
行空間基底を更新し，完全な行空間に対する所属判定を行う．
それでも棄却できない場合のみ，零空間計算に進み，
削減に用いるベクトル $y$ を探索する．

\subsection{実験結果}

提案した早期棄却法を，既存の最速実装である
\texttt{lx2\_dynamic\_repair\_chi\_packed\_aa}
を基に実装した．
実験では，小型 $\chi$ 行列に用いる行数の割合を
$5\%,8\%,10\%,13\%,15\%$ に変化させ，
通常の基底再利用法と比較した．

実験ログ \texttt{slurm-6.out} では，23 個のベンチマーク回路に対して
6 種類の設定を実行した．
各回路で最速の設定を選んだ場合，通常の基底再利用法に対して
幾何平均で約 $2.07$ 倍の高速化が得られた．
ここでの実行時間は，ログ中の \texttt{Execution Time} に対応する
TODD 部分のみの時間である．
合計実行時間は，通常の基底再利用法が $486.6$ 秒であったのに対し，
各回路の最速設定では $172.7$ 秒であった．

設定別の合計実行時間は次の通りである．
\[
\begin{array}{c|c}
\text{手法} & \text{合計実行時間 [s]} \\
\hline
\texttt{aa} & 486.6 \\
\texttt{aa\_sketch\_5} & 480.3 \\
\texttt{aa\_sketch\_8} & 250.9 \\
\texttt{aa\_sketch\_10} & 195.3 \\
\texttt{aa\_sketch\_13} & 177.0 \\
\texttt{aa\_sketch\_15} & 179.8
\end{array}
\]
この結果から，全体としては $13\%$ 設定が最も良好であり，
$15\%$ 設定もほぼ同等の性能を示した．
一方で，$5\%$ 設定では棄却率が低く，
基底更新を十分に削減できないため，効果は限定的であった．

特に，$gf(2^k)$ 乗算回路では大きな効果が確認された．
たとえば，$gf(2^8)$ 乗算では，通常の基底再利用法が
$75.9$ 秒であったのに対し，$13\%$ 設定では $26.2$ 秒となり，
約 $2.7$ 倍の高速化が得られた．
また，$gf(2^9)$ 乗算では，$219.6$ 秒から $73.8$ 秒へ短縮され，
約 $2.8$ 倍の高速化が得られた．
これらの回路では，早期棄却によって候補の大部分が
完全な基底更新の前に除外されるため，基底更新時間が大きく削減された．

\subsection{考察}

今回の結果から，早期棄却法の主な効果は，
完全な $\chi$ 行列に対する基底更新へ進む候補数を大幅に減らす点にあると考えられる．
たとえば，$gf(2^8)$ 乗算の $13\%$ 設定では，
全候補 $24908$ 個のうち $23186$ 個が小型行列によって棄却された．
その結果，完全な基底更新に進む候補は $1706$ 個に抑えられ，
基底更新時間は通常の基底再利用法に比べて大きく減少した．

一方で，小型行列の行数を増やすと棄却率は向上するが，
小型行列自体の構成および所属判定のコストも増加する．
そのため，$15\%$ 設定では $13\%$ 設定より棄却率が高い場合でも，
総実行時間では必ずしも最速にならない場合があった．
このことから，早期棄却では，
棄却率と小型行列の計算コストのバランスが重要である．

また，今回の手法は基底再利用法を土台としているが，
高速化の主因は，基底を大量に局所再利用できたことではなく，
基底更新そのものを行う前に多数の候補を除外できた点にある．
実際，早期棄却後に完全判定へ進んだ候補でも，
局所更新ではなく基底の再構築が必要となる場合は多い．
したがって，本手法は
「基底再利用をさらに効率化する手法」というよりも，
「基底再利用に入る前の候補数を削減する手法」
として位置付けるのが適切である．

\subsection{行サンプリング率ごとの比較}

以下に，$gf(2^7)$ 乗算回路および $gf(2^8)$ 乗算回路に対する
小型 $\chi$ 行列の行サンプリング率ごとの実験結果を示す．
ベースラインには，元の基底再利用法
\texttt{lx2\_dynamic\_repair\_chi\_packed\_aa} を用いた．

\begin{table}[H]
\centering
\caption{$gf(2^7)$ 乗算回路に対する小型 $\chi$ 行列の行サンプリング率ごとの比較}
\label{tab:gf27-sketch-rate}
\begin{tabular}{lrrrrrr}
\hline
手法 & T-count & 合計時間 [s] & 零空間時間 [s] & Sketch 時間 [s] & 零空間+Sketch [s] & 棄却率 [\%] \\
\hline
基底再利用 & $343 \to 164$ & 29.00 & 0.016 & --    & 0.016 & -- \\
$p=5\%$  & $343 \to 164$ & 30.48 & 0.023 & 5.00  & 5.03  & 12.04 \\
$p=8\%$  & $343 \to 164$ & 20.10 & 0.023 & 7.02  & 7.04  & 51.03 \\
$p=10\%$ & $343 \to 164$ & 15.22 & 0.020 & 7.78  & 7.80  & 67.78 \\
$p=13\%$ & $343 \to 164$ & 14.23 & 0.024 & 9.58  & 9.60  & 81.93 \\
$p=15\%$ & $343 \to 164$ & \cellcolor{green}12.97 & 0.022 & 10.18 & 10.20 & 87.33 \\
\hline
\end{tabular}
\end{table}

\begin{table}[H]
\centering
\caption{$gf(2^8)$ 乗算回路に対する小型 $\chi$ 行列の行サンプリング率ごとの比較}
\label{tab:gf28-sketch-rate}
\begin{tabular}{lrrrrrr}
\hline
手法 & T-count & 合計時間 [s] & 零空間時間 [s] & Sketch 時間 [s] & 零空間+Sketch [s] & 棄却率 [\%] \\
\hline
基底再利用 & $448 \to 220$ & 75.89 & 0.022 & --    & 0.022 & -- \\
$p=5\%$  & $448 \to 220$ & 82.67 & 0.031 & 12.21 & 12.24 & 7.68 \\
$p=8\%$  & $448 \to 220$ & 38.88 & 0.025 & 16.63 & 16.65 & 65.37 \\
$p=10\%$ & $448 \to 220$ & 28.66 & 0.028 & 19.26 & 19.29 & 83.68 \\
$p=13\%$ & $448 \to 220$ & \cellcolor{green}26.22 & 0.018 & 22.97 & 22.99 & 93.09 \\
$p=15\%$ & $448 \to 220$ & 27.41 & 0.022 & 25.41 & 25.43 & 95.56 \\
\hline
\end{tabular}
\end{table}

$gf(2^7)$ 乗算回路では，$15\%$ 設定が最速となり，
基底再利用法の $29.00$ 秒から $12.97$ 秒へ短縮された．
一方で，$gf(2^8)$ 乗算回路では，$13\%$ 設定が最速となり，
$75.89$ 秒から $26.22$ 秒へ短縮された．
どちらの回路でも，サンプリング率を上げるほど棄却率は向上するが，
Sketch 時間も増加する．
そのため，最適なサンプリング率は単に棄却率が最大となる設定ではなく，
棄却によって削減される基底更新時間と，Sketch 判定に要する時間の
バランスによって決まる．
